\begin{frame}
    \frametitle{Announcements}

    \begin{alertblock}{Attendance}
        Please make sure to use mobile attendence!!! According to school regulations, attendance is required for at least three-quarters of classes. In other words, if you absent more than quarter of the classes, you will fail the course. 
    \end{alertblock}

    I know mobile attendance sometimes not work, so if you have any issues with it, please send a email to me (anyunpyo@unist.ac.kr) and carbon copy (CC) to the professor (jooyeon.kim@unist.ac.kr). With title as "\textit{[ITP11701] Mobile Attendance Issue - Your Name}". Please include the screenshot of the error message you get when you try to submit mobile attendance. \textbf{Please title as above and CC to the professor, otherwise I might miss your email.}

    Feedback on the mobile attendance system in response to attendance requests will be conducted at \textbf{the end of the semester}.
\end{frame}

\begin{frame}
    \frametitle{Announcements}

    \begin{block}{About using LLMs}
        There are several LLM based services -- GPT, Claude Code, etc. -- that can help you with learning. You should (can) use them for assignments, projects. That said, hands-on coding and self-study would be more valuable than relying on LLM answers.
    \end{block}

    \begin{alertblock}{Individual Work Required}
        All assignments must be completed individually. Copying code from others will result in a negative score. In this case, the assignment in question will be marked as \textcolor{red}{-1} point. Note that lab sessions account for 50\% of the total grade. If there is worry about your submitted code being identical to another student's, adding detailed explanations of the code is one way to address such concerns.
    \end{alertblock}
\end{frame}

\begin{frame}
    \frametitle{Calendar for lab sessions}

    % 10 items with table
    \begin{table}[]
        \begin{tabular}{|c|l|l|}
            \hline
            Week & Date & Topic \\
            \hline
            1 & Mar 4 & No lab session \\
            2 & Mar 11 & \textbf{LAB 0: Review python and jupyter notebooks} \\
            3 & Mar 18 & LAB 1: Linear algebra with numpy and GPU \\
            4 & Mar 25 & LAB 2: Linear regression models \\
            5 & Apr 1 & LAB 3: Logistic regression \\
            6 & Apr 8 & LAB 4: Softmax classifications \\
            7 & Apr 15 & LAB 5: Classification models \\
            8 & Apr 22 & \textcolor{unistemerald}{Midterm break} No lab session \\
            \hline
            9 & Apr 29 & LAB 6: Deep learning from scratch \\
            10 & May 6 & LAB 7: Normalization techniques \\
            11 & May 13 & LAB 8: Convolutional neural networks \\
            12 & May 20 & LAB 9: More on CNN -- ResNet \\
            13 & May 27 & Final project announcement \\
            14 & Jun 3 & \textcolor{red}{Election day} No lab session \\
            15 & Jun 10 & LAB 10: Transformers \\
            16 & Jun 17 & \textcolor{unistemerald}{Final exam week} No lab session \\
            \hline
        \end{tabular}
    \end{table}
\end{frame}

\begin{frame}
  \frametitle{After Lab 0 you can...}
  \begin{itemize}
    \item Set up Python environment on your local machine and Google Colab
    \item Write and execute Python code in Jupyter Notebooks
    \item Use basic Python syntax, data structures and libraries for AI programming
    \item Understanding Jupyter Notebooks features
  \end{itemize}
\end{frame}